\section{Uncertainty structure}\label{sec:8.8}

The local fields of silicon raise its screening factor at small momentum transfer by a quarter (\cref{sec:4.6}), and at a momentum transfer of \SI{7.5}{keV} they lower the screened structure factor by \SI{22}{\percent} to \SI{44}{\percent} between 2 and \SI{10}{eV} (\cref{we:8-lf-nuisance}). A limit computed with the random-phase approximation alone carries that difference whether or not anyone states it. \Cref{sec:1.5} called such errors the theory uncertainty of a limit: they come from approximations, and a longer exposure does not reduce them. \Cref{sec:8.5} gave every approximation of the response a place on the truncation label $\mathfrak{T}$, and with it a next term and, where the axis has a small parameter, a size. What remains is to turn each truncation into a number that a likelihood can use. A theory uncertainty built this way is computed from the physics, and two analyses that use the same truncation share it.

\subsection{The uncertainty of a truncation}

A formula sits at a point $\mathfrak{T}$ of \cref{eq:8.5-label}. Let $\mathfrak{T}_k^+$ be the label one step further along the axis $k$ and nowhere else. On the integer axes the step goes to the next order at which the labeled quantity changes, by the counting rule of \cref{sec:8.5}: $n_g = 0 \to 2$, $n_B = 0 \to 2$, $n_A = 0 \to 2$ and $n_T = 0 \to 1$. On the many-body axis it is the next approximation, from IP to RPA, to RPA + LF, to TDDFT or BSE. The lattice regime is not an order, and its step is the next term of the multiphonon series \cref{eq:5.6-multi}, which is known exactly. The difference of the two responses and its size relative to the first are
\begin{equation}
  \Delta S_k(\vect{q},\omega) = S^{\mathfrak{T}_k^+}(\vect{q},\omega) - S^{\mathfrak{T}}(\vect{q},\omega) , \qquad
  \delta_k(\vect{q},\omega) = \frac{\Delta S_k(\vect{q},\omega)}{S^{\mathfrak{T}}(\vect{q},\omega)} .
  \label{eq:8.8-delta}
\end{equation}
In the likelihood the truncation is represented by a \emph{nuisance parameter} $\vartheta_k$, a parameter that is fitted together with the cross-section without being itself the quantity sought. It interpolates between the two responses. For a single response, such as $S_{ee}$ or the contraction $\sum_{ab}\mathfrak{D}_{ab}S_{ab}$ of one model,
\begin{equation}
  S(\vect{q},\omega;\vartheta) = S^{\mathfrak{T}}(\vect{q},\omega)\prod_k\bigl[1 + \delta_k(\vect{q},\omega)\bigr]^{\vartheta_k} ,
  \label{eq:8.8-M}
\end{equation}
so that $\vartheta_k = 0$ is the formula, $\vartheta_k = 1$ is its next order, and $\vartheta_k$ is given a standard normal constraint in the likelihood. To first order, $S(\vartheta) = S^{\mathfrak{T}} + \sum_k\vartheta_k\ln(1 + \delta_k)\,S^{\mathfrak{T}}$, the familiar additive form. The power form differs from it at second order, and the difference matters: $S^{\mathfrak{T}}(1 + \vartheta_k\delta_k)$ turns negative at $\vartheta_k = -1/\delta_k$, two standard deviations out for $\delta_k = 0.5$, while \cref{eq:8.8-M} is positive for every $\vartheta$.

The parameter $\vartheta_k$ does not vary from bin to bin. A truncation is a property of the calculation, not of the bin, and the same $\vartheta_k$ acts at every $\vect{q}$ and $\omega$, in every bin and on every final-state class $\Lambda$ of \cref{sec:8.7}. It is also shared by every target computed with the same approximation and by every experiment that uses the result. Only the template $\delta_k(\vect{q},\omega)$ varies. This is what distinguishes a nuisance parameter from an error bar attached to each bin. A truncation that lowers the spectrum near the gap and raises it at higher energy is one parameter, and an analysis that treats the two regions as independent errors would let the data fit a shape that no calculation can produce.

When is $\Delta S_k$ computable? The next order may have been computed, it may be estimated from a small parameter, or neither may be possible. In the first case the next order is known, as for the next many-body approximation of \cref{sec:4.6}, the zero-point renormalization of a gap at $n_g = 2$ (\cref{sec:6.3}) or the next term of \cref{eq:5.6-multi}. The template is then a number at every point, and the better result usually becomes the central value, with the uncertainty carried by the step after it. For the many-body approximations that step is again a computed difference, and like every many-body spread it is an estimate, not a bound (\cref{sec:4.9}). In the second case the next order has not been computed, but the axis has a small parameter whose size \cref{tab:8-params} gives. Then $\delta_k$ is the parameter raised to the power of the step, times a coefficient of order one, between a quarter and two in the silicon counting of \cref{sec:5.7}. The uncertainty is then an estimate and is labeled as such. In the third case there is neither: no computed next order and no small parameter. The size of the truncation is then not known, and no width is assigned to it. The pair-channel screening of the phonon-assisted sideband is the example of this chapter. \Cref{sec:8.4} showed that dropping it is an omission, not a choice of convention, and \cref{sec:13.4} gives only its two solvable limits, which bracket it loosely and do not fix its sign. The calculation that would size it has not been done. Until it is, a rate that omits it is reported at its nominal value with the omission named, because a width invented for it would hide the one fact a reader needs.

Approximations off the label carry theory uncertainties too. The exchange--correlation functional, the frozen core and the numerical parameters of a calculation have no next order, and their sizes come from the spreads and convergence tests of \crefrange{chap:23}{chap:26}. They enter the likelihood in the same way, with the spread in place of $\Delta S_k$.

For a response matrix $S_{ab}$ the power form needs care. Raising each element to its own power leaves the cone of positive semidefinite matrices, and $\delta_{k,ab}$ is undefined wherever $S^{\mathfrak{T}}_{ab} = 0$. The matrix analog of \cref{eq:8.8-M} is
\begin{equation}
  S(\vartheta_k) = S^{1/2}\bigl(S^{-1/2}S^+S^{-1/2}\bigr)^{\vartheta_k}S^{1/2} , \qquad S \equiv S^{\mathfrak{T}}(\vect{q},\omega),\ S^+ \equiv S^{\mathfrak{T}_k^+}(\vect{q},\omega) ,
  \label{eq:8.8-geodesic}
\end{equation}
which reduces to \cref{eq:8.8-M} for one row. It equals $S$ at $\vartheta_k = 0$ and $S^+$ at $\vartheta_k = 1$. It is positive semidefinite for every $\vartheta_k \ge 0$, since $S^{-1/2}S^+S^{-1/2}$ is positive semidefinite, so is any nonnegative power of it, and the congruence $X \to S^{1/2}XS^{1/2}$ preserves the property. The inverse square root is taken on the subspace where $S$ is nonzero, which requires $S^+$ to have no weight outside it. A profile explores negative $\vartheta_k$ as well, and a negative power exists only if $S^+$ is positive definite on that subspace: the two truncations must have the same support, $\operatorname{supp}S^+ = \operatorname{supp}S$, and then \cref{eq:8.8-geodesic} is positive semidefinite for every real $\vartheta_k$. Where the next order opens new support, \cref{eq:8.8-open} below applies instead. Several truncations are applied one after another; their order matters only at second order in the $\delta_k$.

\subsection{Scale, tilt and edge shifts}

A truncation uncertainty scales a spectrum, tilts it or moves its edge, and each moves a limit differently.

A \emph{normalization} uncertainty has a template $\delta_k$ that is constant over the energies a rate samples at fixed $\vect{q}$. It scales the spectrum and leaves its shape alone, and by \cref{eq:1.5-limscaling} it moves the limit by the same fraction. It also leaves the weights of the final-state classes unchanged, so it passes to the pair-number spectrum of \cref{eq:8.7-cond} as an overall factor (\cref{sec:8.7}).

A \emph{tilt} has a template that varies across the energies a rate samples. It moves weight between bins, and its effect on a limit depends on which bins the dark-matter spectrum fills, so on the mass and the mediator. A tilt changes the mix of final-state classes as well, and with it the pair-number spectrum.

A \emph{threshold} uncertainty moves an edge,
\begin{equation}
  S(\vect{q},\omega;\vartheta_k) = S^{\mathfrak{T}}(\vect{q},\omega - \vartheta_k\,\delta E_k) ,
  \label{eq:8.8-shift}
\end{equation}
with $\delta E_k$ the shift that the next order produces. A shift acts on the limit through the edge of \cref{sec:1.3-edge}, $\delta\ln m_{\mathrm{min}} = \delta\ln E_{\mathrm{th}}$, and near the edge it moves the limit by a far larger factor than its fractional size, because $\eta(\vmin)$ falls steeply there (\cref{sec:1.5-inputs}).

Some truncations do more than deform the spectrum: their next order puts weight where the truncated formula has none. The exciton lines of silicon below the pair continuum, which dark matter with $\vect{q} \approx \vect{k}_{\mathrm{c}}$ can create directly (\cref{sec:4.9}), are an example, and so is the two-phonon response above the top of the one-phonon band. No power of $1 + \delta_k$ can create weight from zero. A positive piece that the formula drops is an error of known sign, so once computed it belongs in the central value. Its nuisance parameter then acts on its own size,
\begin{equation}
  S(\vect{q},\omega;\vartheta_k) = S^{\mathfrak{T}}(\vect{q},\omega) + \Delta S_k(\vect{q},\omega)\,\bigl(1 + \delta_k'\bigr)^{\vartheta_k} ,
  \label{eq:8.8-open}
\end{equation}
where $\delta_k'$ is the fractional size of the next correction to the new piece. The rate stays positive for every $\vartheta_k$ as long as $\Delta S_k$ is itself a positive spectrum, which holds for new lines and for a sideband inside the off-shell window of \cref{sec:8.6}. Outside the window it fails, since the finite-part terms of \cref{eq:8.6-dyson} are not positive on their own, and there the nuisance parameter acts on the zero-phonon and phonon-dressed sectors together.

\Cref{tab:8-nuisances} gives an example of each type in each sector of \cref{eq:8.6-partition}, for silicon wherever the size is known. Two entries share a nuisance parameter. The quasiparticle residue $z_{\mathrm{qp}}$ lowers the zero-phonon pair weight, and \cref{eq:8.6-zsum} returns exactly that weight to the finite part of the sideband. One parameter acts on both sectors with opposite signs, and their sum does not move. Treated as two independent nuisance parameters, the residue and the sideband would double the uncertainty of a quantity that the identity holds fixed.

\begin{table}[tbp]
  \caption{Types of truncation uncertainty, with one truncation of each type in each sector of the decomposition \cref{eq:8.6-partition}: $\mathcal{P}_{\mathrm{ad}}$ (electrons following the nuclei; phonons and recoils), $\mathcal{P}_0$ (zero-phonon electronic excitations) and $\mathcal{P}_+$ (electronic excitations with phonons). The axis of the truncation label is given in parentheses. Sizes are for silicon, and the mark ``est.'' denotes an estimate from a small parameter of \cref{tab:8-params}. The residue entry in $\mathcal{P}_0$ and the finite-part entry in $\mathcal{P}_+$ carry one nuisance parameter, tied by \cref{eq:8.6-zsum}; so are the two $\mathcal{K}$ entries of $\mathcal{P}_{\mathrm{ad}}$, the two parts of one step of \cref{eq:5.6-multi}.}
  \label{tab:8-nuisances}
  \centering\small\setlength{\tabcolsep}{4pt}
  \begin{tabular}{@{}>{\raggedright\arraybackslash}p{0.14\linewidth}>{\raggedright\arraybackslash}p{0.15\linewidth}>{\raggedright\arraybackslash}p{0.20\linewidth}>{\raggedright\arraybackslash}p{0.20\linewidth}>{\raggedright\arraybackslash}p{0.20\linewidth}@{}}
    \toprule
    Type & Form & $\mathcal{P}_{\mathrm{ad}}$ & $\mathcal{P}_0$ & $\mathcal{P}_+$ \\
    \midrule
    Normali\-zation & \cref{eq:8.8-M}, $\delta_k$ constant in $\omega$ & Debye--Waller factor $\ee^{-2W_\kappa}$ dropped from a one-phonon rate ($\mathcal{K}$): $\delta = -0.061$ at \SI{10}{keV} & residue $z_{\mathrm{qp}}$ of the pair ($n_g = 2$); est.\ of order $\theta_{\mathrm{ep}} \sim 10^{-2}$ & static screening of the vertex ($n_B$): est.\ $(\omega_{\mathrm{ph}}/\Eg)^2 = 3\times10^{-3}$; finite part tied to the residue \\
    Tilt & \cref{eq:8.8-M}, $\delta_k(\omega)$ & anharmonic reshaping of the recoil spectrum at fixed area and mean ($n_A = 2$): \cref{eq:8.8-tilt} & local-field correction, RPA $\to$ RPA + LF ($\mathcal{X}$): $-0.44$ to $+0.25$ across $q$ and $\omega$, \cref{we:8-lf-nuisance} & two more vertices ($n_g = 4$): est.\ $\theta_{\mathrm{ep}} \approx 10^{-2}$ below $E_{\mathrm{dir}}$, rising to 1 near it \\
    Threshold & \cref{eq:8.8-shift} & anharmonic shift of the optical phonon ($n_A = 2$): est.\ $\eta_{\mathrm{anh}}^2\omega_{\mathrm{opt}} = \SI{0.15}{meV}$ & zero-point renormalization of the gap ($n_g = 2$): $\delta E = \SI{-56.2}{meV}$; adiabatic against nonadiabatic \SI{8}{meV} & exciton binding ($\mathcal{X} = \mathrm{IP} \to \mathrm{BSE}$): every phonon-assisted threshold \SI{14.7}{meV} lower \\
    New support & \cref{eq:8.8-open} & two-phonon response above the one-phonon band ($\mathcal{K}$) & exciton lines below the pair continuum at $\vect{q} \approx \vect{k}_{\mathrm{c}}$ ($\mathcal{X} = \mathrm{BSE}$) & phonon-absorption branch ($n_T = 1$): weight $\ee^{-21550}$ at \SI{10}{mK} \\
    \bottomrule
  \end{tabular}
\end{table}

The Debye--Waller entry uses $2W_\kappa = 0.063$ at \SI{10}{keV} (\cref{sec:8.5}): the correct one-phonon rate is lower by $1 - \ee^{-0.063} = \SI{6.1}{\percent}$ than one that sets $\ee^{-2W_\kappa} = 1$, at that momentum and uniformly in $\omega$. The template is exact, and its value grows with $q$. The two-phonon term of the new-support row is the other half of the same step of \cref{eq:5.6-multi}: the weight that $\ee^{-2W_\kappa}$ removes from the one-phonon term goes to it and to the elastic line, and the two entries carry one $\vartheta$. The gap entry is the one in which the central value moves. A clamped-nucleus gap is at $n_g = 0$, and the zero-point renormalization, \SI{-56.2}{meV} in the nonadiabatic form (\cref{sec:6.3}), is its computed next order. Once it is included, the uncertainty that remains is the step after it, of which the \SI{8}{meV} between the adiabatic and nonadiabatic forms and the difference to the measured \SI{-62}{meV} to \SI{-64}{meV} give the size. The anharmonic shift uses the counting $\eta_{\mathrm{anh}}^2\omega_{\mathrm{ph}}$ of \cref{sec:5.7}, with $\eta_{\mathrm{anh}}^2 = 2.4\times10^{-3}$ and $\omega_{\mathrm{opt}} = \SI{64.1}{meV}$. The thermal entry is an uncertainty whose template has been computed and found to vanish: at \SI{10}{mK} the occupations of \cref{sec:8.5} make it smaller than any number a likelihood can resolve.

\subsection{Exact relations as constraints}

The relations of \cref{sec:8.3} hold for the exact response and for every truncation that respects them. A nuisance parameter that breaks one of them describes no possible calculation, and the relations are therefore constraints on the nuisance parameters, not further sources of uncertainty.

\emph{Positivity.} The response matrix is positive semidefinite at every $(\vect{q},\omega)$ (\cref{eq:4.2-positivity}), and $S(\vartheta)$ must be so for every $\vartheta$ that the likelihood explores. The power form \cref{eq:8.8-M}, the matrix form \cref{eq:8.8-geodesic} and the new-support form \cref{eq:8.8-open} with a positive $\Delta S_k$ satisfy this by construction. The additive form $S^{\mathfrak{T}} + \vartheta_k\Delta S_k$ does not, wherever $|\Delta S_k|$ is comparable to $S^{\mathfrak{T}}$, which is common near edges.

\emph{The interference bound.} The cross term of two groups of rows can never exceed twice the geometric mean of their separate rates, \cref{eq:8.3-bound}. A nuisance parameter on the cross term, however, is not free. If a truncation scales the separate rates by $(1 + \delta_A)^{\vartheta}$ and $(1 + \delta_B)^{\vartheta}$, the largest cross term allowed scales by $[(1 + \delta_A)(1 + \delta_B)]^{\vartheta/2}$, and one that grows the cross term faster than the separate rates describes no target. In practice nuisance parameters are assigned to the matrix, through \cref{eq:8.8-geodesic}, and never to an interference term on its own; the bound is then kept automatically. The same holds for the mixed electron--nucleus response of the Migdal effect, where $|S_{e\kappa}|^2 \le S_{ee}S_{\kappa\kappa}$.

\emph{Sum rules.} The first moment of the density response is fixed by the f-sum rule (\cref{eq:4.3-firstmoment}) for any approximation that conserves particle number with local potentials. Every approximation of \cref{sec:4.6} is of that kind: the bubble, the random-phase approximation with or without local fields, and the adiabatic kernel all add static, local terms that leave the tail coefficient of the response unchanged. For these truncations
\begin{equation}
  \int_{-\infty}^{\infty}\dd\omega\;\omega\,\Delta S_k(\vect{q},\omega) = 0 .
  \label{eq:8.8-fsum}
\end{equation}
A many-body uncertainty moves weight in energy but cannot add or remove first moment, and its template must change sign: what it lowers in one range it raises in another. A normalization uncertainty on the whole spectrum would violate \cref{eq:8.8-fsum}, which is why the many-body uncertainties are tilts. The relation is not vacuous for a rate, which samples only energies below $\tfrac12\mchi v_{\mathrm{max}}^2$. Within that window a many-body uncertainty may look like a normalization, and \cref{eq:8.8-fsum} says the compensating weight lies above it. \Cref{sec:4.6} saw the effect in the optical spectrum, where the local fields push weight up. Some truncations do not conserve the f-sum: a cut in the band sum or in $q$, a nonlocal pseudopotential, a scissor shift of the conduction bands. For them the violation of the sum rule measures the error, and the template can be calibrated against it. The partial sum rules of \cref{sec:8.6} act in the same way between sectors, and the residue identity \cref{eq:8.6-zsum} is one of them.

\emph{Detailed balance.} In a warm target the response at negative energies is the Boltzmann-weighted image of the response at positive energies (\cref{eq:2.5-db}), and the dark-matter rates in the two directions are tied by \cref{eq:8.2-db}. A nuisance parameter applied to $\omega > 0$ alone would break the tie. The same $\vartheta_k$ must act on both signs of $\omega$, with the template at $(-\vect{q},-\omega)$ fixed by the one at $(\vect{q},\omega)$.

\emph{Moments of the recoil spectrum.} The self response of a nucleus has a fixed zeroth moment, the norm $\langle\ee^{-\ii\vect{q}\cdot\vect{X}}\ee^{\ii\vect{q}\cdot\vect{X}}\rangle = 1$, which no Hamiltonian changes, and a fixed first moment, the recoil energy $E_R$ of \cref{eq:5.6-sum}. The derivation of \cref{eq:5.6-sum} used only the kinetic energy of the nucleus, so both moments hold on any energy surface, anharmonic ones included. The anharmonic truncation $n_A = 0 \to 2$ can only reshape the recoil spectrum at fixed area and fixed mean. The power form keeps both moments only at $\vartheta = 0$ and 1. A \emph{constrained nuisance parameter} keeps them at every $\vartheta$. With $y = \omega - E_R$ and $\psi = \ln(1 + \delta_A)$, the template of the anharmonic step, it reads
\begin{equation}
  S(\vect{q},\omega;\vartheta) = S^{\mathfrak{T}}(\vect{q},\omega)\,\exp\bigl[\vartheta\,\psi(\omega) + \lambda_0(\vartheta) + \lambda_1(\vartheta)\,y\bigr] ,
  \label{eq:8.8-tilt}
\end{equation}
where $\lambda_0$ and $\lambda_1$ are fixed at each $\vartheta$ by requiring the zeroth and first moments to equal those of $S^{\mathfrak{T}}$. Positivity is manifest. That $\lambda_0$ and $\lambda_1$ exist and are unique follows from one convex function. Write $m_n(\vartheta) = \int\dd\omega\,y^nS(\vect{q},\omega;\vartheta)$ for the moments in the variable $y$, and let $F(\lambda_0,\lambda_1) = \int\dd\omega\,S^{\mathfrak{T}}\exp(\vartheta\psi + \lambda_0 + \lambda_1y) - \lambda_0m_0(0) - \lambda_1m_1(0)$. Its gradient vanishes exactly where $m_0(\vartheta) = m_0(0)$ and $m_1(\vartheta) = m_1(0)$, and its matrix of second derivatives is the moment matrix $\bigl(\begin{smallmatrix}m_0 & m_1\\ m_1 & m_2\end{smallmatrix}\bigr)$ of the deformed spectrum. This is the Hankel matrix of \cref{eq:8.3-hankel}, positive definite for any spectrum with more than one point of support. As a result $F$ is strictly convex and has at most one stationary point, which exists when the prescribed mean lies inside the support, as the free-recoil energy does. Since $S^{\mathfrak{T}_A^+}$ obeys the same two moment conditions, the solution at $\vartheta = 1$ is $\lambda = 0$, and \cref{eq:8.8-tilt} passes through the next order exactly. The second moment is not protected. It contains the mean kinetic energy of the nucleus along $\vect{q}$, which anharmonicity changes, and its change is a separate nuisance parameter with its own computed size. The tilt of \cref{eq:8.8-tilt} is not limited to the nuclear self response: any truncation that must preserve known linear moments uses it, with one $\lambda$ per moment.

\subsection{Nuisance parameters in the likelihood}

A search records counts $N_{\mathrm{obs},j}$ in bins $j$ of the measured quantity, pair number for a charge detector, phonon energy for a calorimeter. Its expected signal in bin $j$ is $N_{\mathrm{sig},j} = \sigmae\,M_{\mathrm{det}}t_{\mathrm{exp}}\,r_j(\mchi;\vartheta)$, with $r_j$ the rate per unit cross-section into the bin, computed from \cref{eq:8.7-handoff} with the response $S(\vartheta)$ in place of $S$. It generalizes the $r(\mchi)$ of \cref{sec:1.5} to one bin and to nonzero nuisance parameters. With a known background $N_{\mathrm{bkg},j}$, the Poisson probabilities of \cref{eq:1.5-poisson} and a standard normal constraint on each $\vartheta_k$, the likelihood of one exposure is
\begin{equation}
  \mathsf{L}(\sigmae,\vartheta) = \prod_j P\bigl(N_{\mathrm{obs},j}\,\big|\,N_{\mathrm{sig},j}(\sigmae,\vartheta) + N_{\mathrm{bkg},j}\bigr)\;\prod_k\ee^{-\vartheta_k^2/2} .
  \label{eq:8.8-like}
\end{equation}
The halo enters $r_j$ at the fixed values of \cref{sec:1.2}, by the convention of \cref{sec:1.5-inputs}, and the detector's own nuisance parameters, efficiency and background, enter in the same way (\cref{sec:21.4}). The nuisance parameters are removed by profiling~\cite[Eq.~(7)]{cowan2011asymptotic}, \cref{eq:C.14-profile} of \cref{app:C.14}: at each $\sigmae$ every $\vartheta_k$ takes the value that maximizes $\mathsf{L}$. To first order in $\vartheta$, a normalization uncertainty $(1 + \delta_k)^{\vartheta_k} \simeq 1 + \vartheta_k\ln(1 + \delta_k)$ is the Gaussian offset of \cref{app:C.14} with the width $\sigma_\vartheta = \ln(1 + \delta_k)$, and the widths add in quadrature as in \cref{eq:C.14-quadrature}.

How much does a theory uncertainty cost a limit? One bin with no events and no background shows the structure. With $N(\vartheta) = N_{\mathrm{nom}}(1 + \delta)^\vartheta$, \cref{eq:8.8-like} gives $-2\ln\mathsf{L} = 2N(\vartheta) + \vartheta^2$ up to a constant. The best fit, at $N_{\mathrm{nom}} = 0$ and $\vartheta = 0$, has the value zero. At fixed $N_{\mathrm{nom}}$ the profile condition $2kN(\hat{\hat\vartheta}) + 2\hat{\hat\vartheta} = 0$, with $k = \ln(1 + \delta)$, gives $\hat{\hat\vartheta} = -kN(\hat{\hat\vartheta})$. The nuisance parameter is pulled toward less signal, since no events were seen, and the profiled statistic is $\mathsf{t}_{\mathrm{p}} = 2u + k^2u^2$, with $u = N(\hat{\hat\vartheta})$. The asymptotic thresholds of \cref{app:C.14} fail at zero counts, and for this illustration we set the threshold to the value $\mathsf{t}^* = 2N_{90} = 4.61$ at which the statistic without nuisance parameter reproduces the exact limit $N_{90} = 2.303$ of \cref{eq:1.5-N90}. Solving $\mathsf{t}_{\mathrm{p}} = \mathsf{t}^*$ for $u$, and $N_{\mathrm{nom}} = u\,\ee^{k^2u}$, the limit grows by the factor
\begin{equation}
  \frac{N_{\mathrm{nom}}^{\mathrm{lim}}}{N_{90}} = \frac{2w\,\ee^{w}}{k^2\mathsf{t}^*} = 1 + \frac{k^2\mathsf{t}^*}{4} + O(k^6) , \qquad w = \bigl(1 + k^2\mathsf{t}^*\bigr)^{1/2} - 1 .
  \label{eq:8.8-cost}
\end{equation}
The cost is of second order in the size of the truncation: \SI{1.0}{\percent} for $\delta = 0.1$, \SI{8}{\percent} for $\delta = 0.3$ and \SI{19}{\percent} for $\delta = 0.5$. Quoting instead the worst case within one standard deviation, $\vartheta = -1$, costs $1 + \delta$, of first order: \SI{10}{\percent}, \SI{30}{\percent} and \SI{50}{\percent}. The two are different statements, a likelihood and an envelope, and a result must say which it reports. The power form matters here as well. A linear form $1 + \vartheta\delta$ can switch the signal off entirely at $\vartheta = -1/\delta$, at the price $\vartheta^2 = 1/\delta^2$ in the likelihood, and for $\delta > (\mathsf{t}^*)^{-1/2} = 0.47$ that price is below the threshold: the profile then excludes nothing at all.

With several bins the data begin to constrain the nuisance parameters themselves. A normalization uncertainty is degenerate with $\sigmae$ and is never constrained by the counts. A tilt or a threshold shift changes the shape of the expected spectrum, and a measured spectrum with enough events can prefer one value of $\vartheta_k$ over another, which an envelope of fixed width cannot express. Since $\vartheta_k$ is shared by every target and experiment that uses the same truncation, a combination of searches fits one value for it, not one per experiment (\cref{sec:21.5}).

\begin{workedexample}{The local-field uncertainty of the random-phase rate in silicon}\label{we:8-lf-nuisance}
A rate computed with the random-phase approximation of the head element alone, \cref{eq:4.6-rpa}, sits at $\mathcal{X} = \mathrm{RPA}$. The next approximation adds the local fields, $\mathcal{X} = \mathrm{RPA + LF}$ of \cref{eq:4.6-rpaLF}, and the template of its uncertainty is $\delta_{\mathcal{X}} = S^{\mathrm{RPA+LF}}/S^{\mathrm{RPA}} - 1$ for the density structure factor. We compute both approximations in the empirical pseudopotential model of silicon used in \cref{we:8-pairs}, with the same bands, the same 51 local-field vectors and the same Lorentzian half-width of \SI{0.10}{eV}, at $q_0 = \SI{3.73}{keV}$ and $2q_0 = \SI{7.46}{keV}$ along $[100]$ and at energies from 0.5 to \SI{100}{eV}. The model gap is \SI{0.82}{eV}, and we quote energies from \SI{2}{eV} up. The quoted values change by 0.01 or less between $16^3$ and $20^3$ points of $\vect{k}$.

At $q_0$ the template is small and has structure. It is $-0.06$ at 2 and \SI{3}{eV}, passes through zero at \SI{4}{eV}, is $-0.03$ at 5 to \SI{7}{eV} and $-0.06$ at \SI{10}{eV}, and reaches its minimum, $-0.15$, near \SI{17.5}{eV}. It crosses zero at \SI{25.6}{eV} and is $+0.02$ at \SI{30}{eV}. At $2q_0$ the local fields matter more: $\delta_{\mathcal{X}} = -0.44$ at \SI{2}{eV}, $-0.30$ at \SI{4}{eV}, $-0.22$ at \SI{10}{eV} and $-0.06$ at \SI{20}{eV}, with the zero near \SI{29}{eV}. At small $q$, finally, the local-field reduction of $\varepsilon_\infty$ alone raises the screening factor $1/\varepsilon_\infty^2$ by $(11.45/10.22)^2 \simeq 1.25$ (\cref{sec:4.6}), so that $\delta_{\mathcal{X}} \approx +0.25$ at the bottom of the spectrum.

The local-field uncertainty is a tilt, in $\omega$ and in $q$. Weighted with the random-phase spectrum between 1.5 and \SI{30}{eV}, its mean is $-0.06$ at $q_0$ and $-0.09$ at $2q_0$. A normalization uncertainty of that size would understate the template near the gap at $2q_0$ by a factor of about five and would have the wrong sign at small $q$. The f-sum rule explains the change of sign. Over 0.5 to \SI{100}{eV} the first moments of the two approximations agree to $1\times10^{-4}$ at $q_0$ and $2\times10^{-4}$ at $2q_0$, the residue that the Lorentzian tails and the finite grid leave, while the zeroth moments fall by \SI{2.1}{\percent} and \SI{3.0}{\percent}. The local fields move weight upward, as \cref{eq:8.8-fsum} demands and as \cref{sec:4.6} found for the optical spectrum, and what they remove below \SI{25}{eV} reappears above it. A particle of mass $\mchi$ deposits at most $\tfrac12\mchi v_{\mathrm{max}}^2 = \SI{3.5}{eV}$ per MeV, so dark matter lighter than about \SI{7}{MeV} samples only energies where $\delta_{\mathcal{X}} \le 0$ at these momenta and never meets the compensating weight.

In the likelihood a single $\vartheta_{\mathcal{X}}$ multiplies $\ln(1 + \delta_{\mathcal{X}})$ at every point. At $\vartheta_{\mathcal{X}} = 1$ the rate is that of the random-phase approximation with local fields. Once that approximation is computed it becomes the central value, and the uncertainty moves to the step to the adiabatic kernel or to the Bethe--Salpeter equation. The adiabatic kernel restores about half of the local-field reduction of $\varepsilon_\infty$ (\cref{sec:4.6}), so its template has the opposite sign at small $q$, and near the gap the Bethe--Salpeter step is the larger one (\cref{we:8-label}).
\end{workedexample}

\section*{Summary}

A slow dark-matter particle sees an electron or a nucleus through the momentum it hands over, their relative velocity and their spins, and nothing else (\cref{sec:8.1}). Every nonrelativistic interaction therefore reduces to five target operators per species: the density, the spin density, the current, the spin current and a kinetic-energy density. The model survives only in the dark-matter-side tensor $\mathfrak{D}_{ab}$, a Hermitian positive semidefinite matrix that weights the response matrix of these operators.

The contraction of $\mathfrak{D}_{ab}$ with $S_{ab}$, integrated over the momentum transfer and the halo, is the master rate formula (\cref{sec:8.2}). It holds at any temperature and for energies of either sign, and its integrand is never negative. Scattering fills a region of the $(q,\omega)$ plane, absorption a point, and the absorption of a fermion a line.

Exact relations constrain every target (\cref{sec:8.3}). Positivity of the response bounds any interference term by twice the geometric mean of the separate rates. The frequency moments of the response form matrices, the odd ones fixed by commutators, the even ones by equal-time correlations, and all of them obey Hankel positivity.

Screening is a property of the coupling, not of the target alone (\cref{sec:8.4}). Only a coupling to the total charge reproduces the loss function at every $(\vect{q},\omega)$, and an absorbed dark photon in silicon at \SI{3}{eV} is suppressed by the factor $|\varepsilon|^2 = 761$ of in-medium mixing.

Every rate is a truncation of one expansion in seven small parameters (\cref{sec:8.5}), among them, in silicon at \SI{10}{mK}, the Born--Oppenheimer parameter 0.0665, the off-shell parameter 0.010 to 0.016 at the edge and the anharmonic parameter 0.049. The truncation label $\mathfrak{T}$ records the order a formula keeps in each and the many-body approximation, lattice regime and Migdal regime it uses.

In a target with a gap the final states split into three sectors (\cref{sec:8.6}): electrons following the nuclei, electronic excitations without phonons and electronic excitations with phonons. The split is unique to the order of the off-shell parameter inside the off-shell window.

The final-state-resolved response sums to the full response (\cref{sec:8.7}). At $q_0$ and \SI{10}{eV} in silicon, dark matter leaves the electron with \SI{4.8}{eV} and the hole with \SI{4.4}{eV} on average, while a photon of the same energy gives \SI{5.7}{eV} and \SI{3.5}{eV}. The cascade must therefore start from the pair distribution at the momentum transfer that made it.

Each truncation becomes a nuisance parameter whose size is the computed step to its next order (\cref{sec:8.8}). These parameters act as normalizations, tilts or threshold shifts, or open new support, and the exact relations constrain them. In the likelihood of one exposure a profiled nuisance parameter costs a limit at second order in its size, an envelope at first order. In silicon the step from the random-phase approximation to local fields is such a tilt, from $-0.44$ to $+0.25$ across momentum and energy.
